\wprowadzenie

Przewodnik obejmuje cały cykl seminarium: od pierwszego uruchomienia wzoru, przez budowę pakietu i pisanie kolejnych rozdziałów, po kontrolę formalną przed oddaniem pracy.

Praca dyplomowa na kierunku elektroniczne przetwarzanie informacji składa się z dwóch ściśle powiązanych części: projektu dyplomowego oraz pracy licencjackiej, która ten projekt opisuje. Zaliczenie seminarium następuje dopiero po przyjęciu projektu \parencite{standardyEPI}. Twoim projektem jest pakiet języka R implementujący algorytm z artykułu metodycznego wraz z serwisem, który go udostępnia.

\section{Dla kogo i po co}

Przewodnik zakłada, że znasz podstawy statystyki i programowania w R \parencite{rcore}, a nie znasz LaTeX-a ani inżynierii oprogramowania. Wszystkiego, co potrzebne, nauczysz się tutaj od zera.

Dokument, który czytasz, jest jednocześnie instrukcją i przykładem. Złożono go tą samą klasą, z której korzystasz przy pisaniu pracy, i tym samym stylem bibliograficznym. Każdy element opisany w rozdziale pierwszym -- tabela z tytułem nad nią, wzór z etykietą, listing kodu, cytowanie w nawiasie, spis ilustracji -- występuje dalej w praktyce. Kiedy nie pamiętasz zapisu, zajrzyj do źródła tego przewodnika: leży w tym samym repozytorium.

\section{Jak czytać}

Rozdziały pierwszy, drugi i trzeci dotyczą narzędzi. Przeczytaj je na początku seminarium, zanim napiszesz pierwsze zdanie pracy. Rozdział pierwszy wystarczy przerobić do połowy, żeby zacząć -- reszta przyda się później.

Rozdziały czwarty, piąty i szósty dotyczą pisania. Do rozdziału czwartego wracaj przy każdej kolejnej części pracy. Rozdział piąty przeczytaj raz, na samym początku, bo pytanie badawcze stawia się przed pisaniem, a nie po nim. Rozdział szósty służy jako lista kontrolna przed oddaniem.

Rozdziały siódmy i ósmy dotyczą pracy z narzędziem wspomagającym pisanie kodu. Siódmy opisuje metodę, ósmy granice dozwolonego użycia i sposób jego dokumentowania. Oba przeczytaj \wyroznienie{zanim wydasz pierwsze polecenie}, a nie wtedy, gdy pakiet będzie już gotowy: część wymagań dotyczy rzeczy, których nie da się odtworzyć po fakcie.

\begin{tabelaepi}{Zawartość przewodnika}{tab:zawartosc}
\begin{tabular}{@{}cp{6.4cm}p{6.2cm}@{}}
\toprule
Rozdział & Temat & Kiedy czytać \\
\midrule
1 & LaTeX & na początku, pierwsze dwa poziomy \\
2 & Bibliografia i Zotero & przed pierwszą lekturą \\
3 & Overleaf i współpraca & przy zakładaniu projektu \\
4 & Jak napisać pracę & przy pisaniu każdej części \\
5 & Warsztat akademicki & raz, na samym początku \\
6 & Wymogi formalne & przed oddaniem \\
7 & Praca z agentem programistycznym & przed pierwszym poleceniem \\
8 & Jawność i odpowiedzialność & na początku i przed oddaniem \\
\bottomrule
\end{tabular}
\end{tabelaepi}

\section{Dokumenty rozstrzygające}

Wymagania wobec pracy określają dwa dokumenty Instytutu. \wyroznienie{Standardy prac dyplomowych EPI} \parencite{standardyEPI} rozstrzygają o układzie pracy i o wymogach wobec projektu dyplomowego. \wyroznienie{Instrukcja ISI} \parencite{instrukcjaISI} rozstrzyga o składzie, przypisach i opisach bibliograficznych. Tam, gdzie dokumenty się rozchodzą, w przewodniku wskazany jest ten, który obowiązuje.

Wymagania te są we wzorze zrealizowane automatycznie. Nie zmieniaj ustawień składu -- rozbieżność z wymogami obciąża pracę, a nie wzór.

\section{Dwie ścieżki pracy}\label{sec:sciezki}

Pracę można pisać na dwa sposoby i przewodnik opisuje oba.

\wyroznienie{Ścieżka przeglądarkowa.} Piszesz na uczelnianej instalacji Overleaf.
Nie instalujesz niczego, kompilacja dzieje się na serwerze, promotor widzi tę samą
wersję co Ty. Jest to ścieżka domyślna i wystarczająca do napisania całej pracy.

\wyroznienie{Ścieżka lokalna.} Piszesz u siebie, w Positronie, mając zainstalowaną
dystrybucję \TeX-a i R. Zyskujesz to, że praca, pakiet i wykresy są w jednym miejscu,
oraz dostęp do skryptów kontrolnych, których na Overleaf uruchomić się nie da.

Różnica jest istotna w kilku miejscach: przy kompilacji, przy czyszczeniu plików
pomocniczych i przy kontroli objętości. Wszędzie tam, gdzie postępowanie się rozchodzi,
akapity są oznaczone słowami \wyroznienie{Lokalnie} i \wyroznienie{Na Overleaf}.
Tam, gdzie takiego oznaczenia nie ma, opis dotyczy obu ścieżek.

Najczęstszy układ w praktyce jest mieszany: pakiet i wykresy powstają lokalnie,
tekst pracy na Overleaf. Wtedy stosujesz oznaczenie właściwe dla czynności, a nie
dla całej pracy.
